% ======================================================================
% sections/methods.tex
%
% Methods scaffold for the paper:
%   Triple Humanoid 24/7 Adverse Event Oncology Trial Response Team:
%   4-Site Rotation
%
% This file is a draft scaffold. A future Claude Code Opus 4.7 1M Max
% session reads the bracketed instruction blocks below, then reads the
% files those instructions name, then expands the scaffold into the
% finished methods section. Target length is roughly 18 to 22 pages of
% the 70 plus page final paper.
% ======================================================================

\begin{sectionaccent}

[METHODS INSTRUCTIONS FOR FUTURE CLAUDE CODE OPUS 4.7 1M MAX SESSION.
Write the methods section as 9 subsections inside this single
\texttt{\textbackslash section\{Methods\}} block. Read every file named
below before writing the subsection that uses it. Do not invent
implementation details. Use the file names verbatim in the prose so a
reader can navigate to the source code.

Subsection 1, Code Instruction Layer. Read:
\begin{itemize}
  \item \texttt{demo-projects/07-humanoid/paper/instructions/README.md}
        (whole file).
  \item Every \texttt{README.md} in the 10 numbered subdirectories of
        \texttt{demo-projects/07-humanoid/paper/instructions/}:
        \texttt{00-project-overview/}, \texttt{01-commit-roadmap/},
        \texttt{02-config-instructions/},
        \texttt{03-schemas-instructions/},
        \texttt{04-robotic-instructions/},
        \texttt{05-cartesian-instructions/},
        \texttt{06-iteration-instructions/},
        \texttt{07-llm-planner-instructions/},
        \texttt{08-comparison-instructions/}, and
        \texttt{09-runtime-instructions/}. The
        \texttt{10-repository-update-instructions/} directory is
        named in the instructions README repository structure block
        but is not yet present on disk; mention this gap explicitly
        as a documented future commit.
\end{itemize}

Subsection 2, Code Generation Layer And The 27 \texttt{src/} Files.
Read every file in
\texttt{demo-projects/07-humanoid/paper/codegen/src/} and discuss the
function of each. Group the discussion by function rather than by
filename:
\begin{itemize}
  \item Sensor to Cartesian pipeline: \texttt{sensor\_to\_xyz.py},
        \texttt{kinematics.yaml}, \texttt{cartesian\_planner.py},
        \texttt{xyz\_safety.py}. Discuss how raw sensor frames become
        x, y, z coordinate goals subject to per arm cumulative force
        budget caps of 15 N during chest compressions and 5 N
        otherwise, and to an inter robot minimum distance of 1.2 m at
        rest and 0.4 m during shared task handoff, and a cumulative
        cross robot force cap of 22 N during 3 robot patient
        transfer.
  \item Per site dispatch: \texttt{h2\_dispatcher.py}. This file
        exhibits per site H2 robot dispatch proficiency and is the
        single point where role assignment becomes a robot command.
        Call out the file by name in the prose.
  \item Swarm camaraderie: \texttt{swarm\_coordinator.py} at 191 lines.
        This file implements the camaraderie logic that animates the
        3 robot Unitree H2 EDU camarade swarm at one site. It is used
        by the LLM planner to author broadcasts and by the in-robot
        loop to update local behavior. Discuss that 191 lines is a
        compact length for non trivial 3 way coordination logic; the
        compactness is a direct consequence of using the per-site
        broadcaster pattern and the shared state tracker rather than
        a peer to peer consensus protocol.
  \item LLM planner and broadcast: \texttt{llm\_planner.py},
        \texttt{broadcaster.py}, \texttt{prompt\_frozen.md},
        \texttt{comms\_intellectual.py}, \texttt{comms\_physical.py},
        \texttt{cross\_site\_bus.py}, \texttt{peer\_state\_tracker.py}.
        Note the 1 Hz broadcast cadence and the redundant failover
        within the site.
  \item CTCAE grading and FDA submission:
        \texttt{ctcae\_grader.py},
        \texttt{fda\_rtct\_submitter.py} at 67 lines. The submitter
        submits CTCAE grade 2 or higher AE records to the FDA RTCT
        pilot intake API within 1 hour per the HR 9505 SLA, emits
        records matching the \texttt{fda\_rtct\_submission} schema,
        and uses ed25519 signing in production with a placeholder in
        simulation. Discuss that 67 lines is the right scale for a
        signed JSON submitter and that the small surface area is a
        regulatory feature, not a limitation.
  \item Physician escalation and error checking:
        \texttt{physician\_escalation.py},
        \texttt{check\_errors.py}. The error checker provides LLM
        credibility to FDA and other regulators by trapping schema
        and physics rule violations before broadcast.
  \item Site runtime and iteration sweep:
        \texttt{site\_runtime.py}, \texttt{iterate.py},
        \texttt{week\_runner.py}, \texttt{ingest.py},
        \texttt{compute.py}.
  \item Comparison and figures: \texttt{compare\_agent.py},
        \texttt{figures\_gen.py}.
  \item Native code paths: \texttt{robot\_loop.cpp} for the in-robot
        10 Hz loop, \texttt{runner.rs} and \texttt{runner\_main.rs}
        for the Rust runner.
\end{itemize}
List the function of every file. The 27 files in
\texttt{codegen/src/} taken together represent a processing feat
because they cover dispatch, planning, kinematics, comms, grading,
submission, error checking, iteration, runtime, comparison, figure
generation, native robot loop, and Rust runner inside one tree.

Subsection 3, Schemas And Configs. Read every file in
\texttt{demo-projects/07-humanoid/paper/codegen/schemas/} and
\texttt{demo-projects/07-humanoid/paper/codegen/config/}.

Subsection 4, Cartesian Coordinate Surface. Read
\texttt{demo-projects/07-humanoid/paper/codegen/src/kinematics.yaml},
\texttt{cartesian\_planner.py}, \texttt{sensor\_to\_xyz.py}, and
\texttt{xyz\_safety.py} again with a focus on the x, y, z command
representation. Discuss that the Cartesian surface is the contract
between the LLM and the robot; the LLM does not author joint angles.

Subsection 5, Camaraderie Behavior. Read the 191 line
\texttt{swarm\_coordinator.py} and the ASCII diagram
\texttt{demo-projects/07-humanoid/paper/codegen/diagrams/swarm\_dance.txt}.
Include the diagram verbatim inside an \texttt{asciibox} environment
as Figure 3.

Subsection 6, Iteration Sweep. Read
\texttt{demo-projects/07-humanoid/paper/codegen/data/iterations/index.jsonl}.
Quote the 32 iteration count across 4 trial sites and many columns as
direct evidence that Claude Code can create project code effectively.
Also read \texttt{demo-projects/07-humanoid/paper/execution/diagrams/03\_iteration\_sweep\_funnel.txt}
and include it verbatim inside an \texttt{asciibox} environment as
Figure 4.

Subsection 7, Tests And Error Checks. Read
\texttt{demo-projects/07-humanoid/paper/codegen/src/check\_errors.py}
and \texttt{demo-projects/07-humanoid/paper/codegen/tests/test\_xyz\_safety.py}.
Both files provide LLM credibility in generating tests and error
checks that the FDA and other regulators would expect in any human
factors review.

Subsection 8, Runtime Targets. Read every file in
\texttt{demo-projects/07-humanoid/paper/codegen/docker/} and the
\texttt{docker-compose.yml} at the codegen tree root. Cover MacOS,
Windows, and Linux runtime targets per the
\texttt{09-runtime-instructions/} directory.

Subsection 9, Execution. Read
\texttt{demo-projects/07-humanoid/paper/execution/README.md} and the
log files in
\texttt{demo-projects/07-humanoid/paper/execution/logs/}
(numbered 01 through 23). Walk through the schema ingest, smoke test,
cargo build, Rust runner, validation pass log progression so a
regulator can audit the run.

Required ASCII diagrams to include verbatim in this section:
\begin{itemize}
  \item \texttt{paper/codegen/diagrams/swarm\_dance.txt} as Figure 3.
  \item \texttt{paper/codegen/diagrams/ae\_response\_flow.txt} as
        Figure 5, placed at the end of subsection 9.
  \item \texttt{paper/execution/diagrams/01\_module\_dependency.txt}
        as Figure 6.
  \item \texttt{paper/execution/diagrams/03\_iteration\_sweep\_funnel.txt}
        as Figure 4 (subsection 6).
\end{itemize}

Required real-life feasibility insight to include in this section:
``Every Cartesian goal in \texttt{cartesian\_planner.py} is bounded by
the force budget caps and the inter robot minimum distance constants
named in the instructions README. These bounds are conservative
relative to what the H2 EDU spec sheet allows. The implication is
that a future regulator review can audit the safety envelope from
the constants alone, without re-running the simulation.''

Do not exceed 22 pages. Use single dashes only. Use \S\ for the
section sign. Avoid stranding 1 or 2 word lines on their own.]

The methods section will walk through the instruction layer, the 27
source files in \texttt{codegen/src/}, the schemas and configs, the
Cartesian coordinate command surface, the camaraderie behavior, the
32 iteration sweep, the test and error check files, the runtime
targets, and the execution log progression.

\sectiontable{Methods planning matrix}{%
Instructions & M1 & 1 & instructions tree & 10 subdir walkthrough & 1 to 2 pages & Names every subdir & Ready \\
Source files & M2 & 2 & codegen/src 27 files & Per file function & 4 to 5 pages & swarm\_coordinator at 191 & Ready \\
Schemas & M3 & 3 & schemas plus config & Schema by schema & 2 pages & Every JSON named & Ready \\
Cartesian & M4 & 4 & kinematics plus planner & x, y, z contract & 2 pages & Force caps named & Ready \\
Camaraderie & M5 & 5 & swarm\_coordinator.py & 191 line walkthrough & 2 pages & Figure 3 ASCII & Ready \\
Iteration & M6 & 6 & iterations/index.jsonl & 32 sweep summary & 1 to 2 pages & 32 across 4 sites & Ready \\
Tests & M7 & 7 & check\_errors plus tests & FDA credibility & 1 to 2 pages & Both files named & Ready \\
Runtime & M8 & 8 & docker plus compose & 3 OS targets & 1 to 2 pages & MacOS Win Linux & Ready \\
Execution & M9 & 9 & execution logs 01 to 23 & Run progression & 2 to 3 pages & Audit trail & Ready \\
}

\end{sectionaccent}
